\documentclass[12pt]{article}
\usepackage{notes-project}

\begin{document}

\begin{center}
  {\LARGE\bfseries Lecture 3: Limit of Sequence}\\[0.25em]
  {\large Mathematical Analysis}\\[0.15em]
\end{center}

\section{Convergent Sequences} 

\begin{dfn}[Convergent Sequence]{convergent-sequence}
  A sequence $a: \NNOne \to F$, where $F$ is a \xref[ordered field]{lec01::def:ordered-field}, is said to be \emph{convergent} if
  \[
    \exists \ell \in F,\, \forall \varepsilon \in F,\, \varepsilon > 0,\, \exists n \in \NNOne,\, \forall m \in \NNOne,\, m \geq n,\, a_m \in \neighbourhood{\varepsilon}{\ell}.
  \]
\end{dfn}

\begin{rem}
The intuition of convergence is that we can find a value $\ell$ such that we can propose any $\epsilon > 0$, such that there exists an index $n \in \NNOne$, for all $m \in \NNOne$ where $m > n$, 
\end{rem}

\begin{rem}[Convergence Over a General Ordered Field]{convergence-ordered-field}
Results in this lecture are stated for a general \xref[ordered field]{lec01::def:ordered-field} $F$ by default; theorems that require a \xref[complete ordered field]{lec01::def:complete-ordered-field} will say so explicitly in their hypotheses. Recall that $\QQ$ is an ordered field and $\RR$ is a complete ordered field.
\end{rem}

\section{Limit of a Sequence} 

\section{Bounded Sequences} 

\section{Monotone Convergence Theorem} 

\end{document}
